\chapter{Postulates and the language of quantum mechanics}
\label{ch:postulates}

\section*{Chapter overview}

Quantum mechanics is the framework that replaces classical mechanics
when the systems we describe are small enough or cold enough that
their phase-space volume becomes comparable to $\hbar$. Before doing
any physics with it, we need a vocabulary: the Hilbert space in which
states live, the algebra of operators that represent observables, and
the rules that govern measurement and time evolution. This chapter
collects exactly that vocabulary. The development is concise and
deliberately a-physical; intuition for what these objects mean will be
built in Chapters~\ref{ch:model-systems}--\ref{ch:perturbations} once
we apply the toolbox to concrete systems.

% --------------------------------------------------------------
\section{States and Hilbert space}
\label{sec:states}

\begin{phenomenon}
Single-photon and single-electron interference experiments
demonstrate that quantum systems do not, in general, possess a
definite value of every observable prior to measurement. The full
list of possible outcomes for any measurement, together with the
probability of each, is encoded in a single abstract object that we
call the \emph{state} of the system.
\end{phenomenon}

\begin{postulate}[State space]
\label{post:states}
The state of an isolated quantum system is described by a unit
vector $\ket{\psi}\in\hilbert$ in a complex separable Hilbert space.
Two unit vectors that differ only by an overall phase, $\ket{\psi}$
and $e^{i\alpha}\ket{\psi}$, are \emph{physically equivalent}: no
measurement can distinguish them, because every observable
prediction depends only on $|\braket{\phi}{\psi}|^2$, in which the
global phase $e^{i\alpha}$ cancels (cf.\ the Born rule,
\cref{post:born}).
\end{postulate}

\begin{definition}[Hilbert space]
\label{def:hilbert}
A \emph{Hilbert space} $\hilbert$ is a complex vector space equipped
with an inner product
$\braket{\cdot}{\cdot}:\hilbert\times\hilbert\to\C$
(\cref{def:inner-product}: sesquilinear, conjugate-symmetric, and
positive definite; \cref{def:sesquilinear,def:conj-sym,def:positive-definite})
that is complete under the induced norm
$\|\psi\|=\sqrt{\braket{\psi}{\psi}}$ in the sense of
\cref{def:cauchy}. The Hilbert spaces relevant for these notes are
all \emph{separable} (\cref{def:separable}), i.e.\ they admit a
countable orthonormal basis.
\end{definition}

In Dirac (bra--ket) notation we write vectors as kets $\ket{\psi}$
and elements of the dual space (\cref{def:dual-space}) as bras
$\bra{\psi}$; the action of a bra on a ket is the inner product,
$\braket{\phi}{\psi}\in\C$. A countable orthonormal basis
$\{\ket{n}\}$ of $\hilbert$ (\cref{def:separable}) gives the
\emph{resolution of the identity}
\begin{equation}
\label{eq:resolution-identity}
\id \;=\; \sum_{n} \ket{n}\bra{n},
\end{equation}
which we will use ubiquitously to insert complete sets of states into
expressions.

\begin{example}[The simplest non-trivial Hilbert space]
The smallest Hilbert space that is not just $\C$ itself is
$\hilbert=\C^2$, with orthonormal basis $\{\ket{0},\ket{1}\}$.
Every unit vector in $\hilbert$ can be written, up to an overall
phase, as
\begin{equation}
\ket{\psi} = \cos(\theta/2)\,\ket{0} + e^{i\varphi}\sin(\theta/2)\,\ket{1},
\end{equation}
parametrised by a point $(\theta,\varphi)$ on the unit sphere $S^2$.
This is the Hilbert space of a two-level system (introduced as a
quantum-mechanical model in \cref{sec:tls}); the geometric picture
of the unit sphere will reappear there as the Bloch sphere
(\cref{sec:bloch}).
\end{example}

\paragraph{Transition.}
States by themselves are inert: they only become physical objects
once we say what we can \emph{measure} on them. Measurements are
mediated by operators, to which we now turn.

% --------------------------------------------------------------
\section{Observables and operators}
\label{sec:operators}

\begin{phenomenon}
Some physical quantities --- energy, position, spin component along
a given axis --- can in principle be measured, while others cannot
be measured simultaneously without disturbance. Both facts must be
encoded in the mathematical object that represents a measurable
quantity.
\end{phenomenon}

\begin{postulate}[Observables]
\label{post:observables}
Every physical observable $A$ is represented by a self-adjoint
(Hermitian) linear operator $\op{A}$ on $\hilbert$. The possible
outcomes of a measurement of $A$ are the eigenvalues of $\op{A}$.
\end{postulate}

\begin{definition}[Hermitian, unitary, projector]
\label{def:operator-types}
Let $\op{A}$ be a linear operator on $\hilbert$
(\cref{def:linear-op}), and $\op{A}^\dagger$ its adjoint
(\cref{def:adjoint}). We call $\op{A}$
\begin{itemize}
  \item \emph{Hermitian} (or self-adjoint) if $\op{A}^\dagger=\op{A}$;
  \item \emph{unitary} if $\op{A}^\dagger\op{A}=\op{A}\op{A}^\dagger=\id$;
  \item a \emph{projector} if $\op{A}^2=\op{A}$ and $\op{A}^\dagger=\op{A}$.
\end{itemize}
\end{definition}

\begin{theorem}[Spectral theorem, finite-dimensional version]
\label{thm:spectral}
Any Hermitian operator $\op{A}$ on a finite-dimensional Hilbert
space admits an orthonormal eigenbasis. Equivalently, there exist
real numbers $\{a_n\}$ and a complete set of orthogonal projectors
$\{\op{P}_n\}$ with $\op{P}_n\op{P}_m=\delta_{nm}\op{P}_n$ and
$\sum_n\op{P}_n=\id$ such that
\begin{equation}
\label{eq:spectral-decomp}
\op{A} \;=\; \sum_n a_n\,\op{P}_n.
\end{equation}
\end{theorem}

\begin{proof}[Sketch]
For finite-dimensional $\hilbert$ the proof reduces to the linear
algebra fact that any Hermitian matrix is unitarily diagonalisable.
For infinite-dimensional $\hilbert$ a more careful spectral measure
is required; we will not need that level of generality here.
\end{proof}

\begin{remark}
Unitaries play a dual role: as operators that conserve the inner
product,
$\braket{\hat{U}\phi}{\hat{U}\psi}=\braket{\phi}{\psi}$, and
therefore preserve probabilities; and as the operators that implement
reversible time evolution (\cref{sec:time-evolution}).
\end{remark}

\paragraph{Transition.}
Equipped with states and operators, we can now state how the two
combine to give an actual measurement outcome.

% --------------------------------------------------------------
\section{Measurement and the Born rule}
\label{sec:measurement}

\begin{phenomenon}
A measurement of a quantum system in general yields one of several
possible outcomes, with probabilities determined by the state. After
the measurement, the system's state is updated to be consistent with
the outcome that was observed.
\end{phenomenon}

\begin{postulate}[Born rule and projective measurement]
\label{post:born}
Let $\ket{\psi}\in\hilbert$ and let $\op{A}=\sum_n a_n\op{P}_n$ be
an observable. A measurement of $\op{A}$ yields the outcome $a_n$
with probability
\begin{equation}
\label{eq:born}
p(a_n) \;=\; \bra{\psi}\op{P}_n\ket{\psi},
\end{equation}
and immediately after the measurement the state is updated to
\begin{equation}
\label{eq:projection}
\ket{\psi}\;\longmapsto\; \frac{\op{P}_n\ket{\psi}}{\sqrt{p(a_n)}}.
\end{equation}
\end{postulate}

\begin{example}[Measuring $\hat\sigma_z$ on a qubit]
Take $\op{A}=\hat\sigma_z$ with eigendecomposition
$\hat\sigma_z=(+1)\ket{0}\bra{0}+(-1)\ket{1}\bra{1}$. For
$\ket{\psi}=\cos(\theta/2)\ket{0}+e^{i\varphi}\sin(\theta/2)\ket{1}$
the Born rule gives $p(0)=\cos^2(\theta/2)$ and
$p(1)=\sin^2(\theta/2)$, independent of $\varphi$.
\end{example}

The expectation value of $\op{A}$ in state $\ket{\psi}$ is the
classical average of the outcomes weighted by their probabilities,
\begin{equation}
\label{eq:expectation}
\expval{\op{A}}_\psi \;=\; \sum_n a_n\,p(a_n)
\;=\; \bra{\psi}\op{A}\ket{\psi}.
\end{equation}

\paragraph{Transition.}
The fact that two observables share or fail to share an eigenbasis
controls whether they can be measured jointly. This is what
commutators encode.

% --------------------------------------------------------------
\section{Commutators and the uncertainty principle}
\label{sec:uncertainty}

\begin{definition}[Commutator and anticommutator]
For operators $\op{A},\op{B}$ define
\begin{equation}
[\op{A},\op{B}] = \op{A}\op{B}-\op{B}\op{A},
\qquad
\{\op{A},\op{B}\} = \op{A}\op{B}+\op{B}\op{A}.
\end{equation}
\end{definition}

Two Hermitian operators share an orthonormal eigenbasis if and only
if $[\op{A},\op{B}]=0$. When $[\op{A},\op{B}]\neq0$ the two
observables cannot be sharply specified at the same time, in a sense
made quantitative by the following theorem.

\begin{theorem}[Robertson uncertainty relation]
\label{thm:robertson}
For any state $\ket{\psi}$ and Hermitian operators $\op{A},\op{B}$,
\begin{equation}
\label{eq:robertson}
\Delta A\,\Delta B \;\geq\; \tfrac{1}{2}\,
\bigl|\expval{[\op{A},\op{B}]}_\psi\bigr|,
\end{equation}
where $\Delta A=\sqrt{\expval{\op{A}^2}-\expval{\op{A}}^2}$.
\end{theorem}

\begin{proof}
Define the shifted operators
$\delta\op{A}=\op{A}-\expval{\op{A}}\,\id$ and similarly
$\delta\op{B}$, and the auxiliary state
$\ket{f}=\delta\op{A}\ket{\psi}$,
$\ket{g}=\delta\op{B}\ket{\psi}$. The Cauchy--Schwarz inequality
gives
\begin{equation}
\bigl|\braket{f}{g}\bigr|^2 \;\leq\; \braket{f}{f}\,\braket{g}{g}
\;=\; (\Delta A)^2(\Delta B)^2.
\end{equation}
Decomposing
$\braket{f}{g}=\expval{\delta\op{A}\,\delta\op{B}}=
\tfrac12\expval{\{\delta\op{A},\delta\op{B}\}}+\tfrac12\expval{[\op{A},\op{B}]}$,
the first term is real and the second imaginary, so
$|\braket{f}{g}|^2\geq\tfrac14|\expval{[\op{A},\op{B}]}|^2$.
Combining the two inequalities yields \cref{eq:robertson}.
\end{proof}

\begin{example}[Position and momentum]
With $[\op{x},\op{p}]=i\hbar\,\id$ the Robertson relation reduces to
the familiar Heisenberg inequality $\Delta x\,\Delta p\geq\hbar/2$.
\end{example}

\paragraph{Transition.}
Static states and a probability rule are not yet a theory of motion.
We close the chapter by giving the equation that drives the state
forward in time.

% --------------------------------------------------------------
\section{Time evolution}
\label{sec:time-evolution}

\begin{postulate}[Schr\"odinger equation]
\label{post:schrodinger}
The state of an isolated quantum system evolves continuously in
time according to
\begin{equation}
\label{eq:schrodinger}
i\hbar\,\dv{t}\ket{\psi(t)} \;=\; \op{H}(t)\,\ket{\psi(t)},
\end{equation}
where $\op{H}(t)$ is the Hermitian Hamiltonian operator.
\end{postulate}

For a time-independent Hamiltonian the formal solution is
$\ket{\psi(t)}=\op{U}(t)\ket{\psi(0)}$ with the unitary
$\op{U}(t)=e^{-i\op{H}t/\hbar}$. For time-dependent Hamiltonians
the solution requires the time-ordered exponential
\begin{equation}
\op{U}(t,t_0) = \mathcal{T}\exp\!\Bigl[-\tfrac{i}{\hbar}\!\int_{t_0}^t\!\op{H}(t')\,\dd t'\Bigr],
\end{equation}
which we will need in \cref{ch:perturbations}.

\subsection{Schr\"odinger, Heisenberg and interaction pictures}
\label{sec:pictures}

\paragraph{The freedom of choice.}
Every physical prediction is an expectation value of the form
$\bra{\psi(t)}\op{O}\ket{\psi(t)}$. Inserting any unitary
$\op{W}(t)$ in the middle ($\op{W}^\dagger\op{W}=\id$) leaves the
prediction unchanged:
\begin{equation}
\bra{\psi(t)}\op{O}\ket{\psi(t)}
\;=\; \bra{\psi(t)}\op{W}^\dagger\,\bigl(\op{W}\op{O}\op{W}^\dagger\bigr)\,\op{W}\ket{\psi(t)}.
\end{equation}
By choosing $\op{W}$ we can shift time dependence between the state
and the operator without changing any measurable quantity. The
three standard choices are:
\begin{itemize}
  \item \emph{Schr\"odinger picture}: $\op{W}=\id$. States carry
        all the time dependence ($\ket{\psi(t)}_S$); operators are
        time-independent ($\op{O}_S$). This is the picture in
        which we stated the Schr\"odinger equation
        \cref{eq:schrodinger}.
  \item \emph{Heisenberg picture}: $\op{W}=\op{U}^\dagger(t)$ with
        $\op{U}(t)=e^{-i\op{H}t/\hbar}$. The unitary
        $\op{W}=\op{U}^\dagger$ exactly undoes the Schr\"odinger
        evolution of the state, so states become time-independent
        and operators inherit the dynamics.
  \item \emph{Interaction picture}: $\op{W}=\op{U}_0^\dagger(t)$
        with $\op{U}_0(t)=e^{-i\op{H}_0 t/\hbar}$ for some
        \emph{reference} Hamiltonian $\op{H}_0$ (a piece of the
        full $\op{H}=\op{H}_0+\op{V}(t)$ that we want to ``rotate
        out''). States and operators both evolve, but only under
        the parts of $\op{H}$ that the reference does not capture.
\end{itemize}

\begin{statement}[Picture transformations]
\label{stmt:pictures}
Let $\op{H}=\op{H}_0+\op{V}(t)$,
$\op{U}(t)=e^{-i\op{H}t/\hbar}$ (full evolution) and
$\op{U}_0(t)=e^{-i\op{H}_0 t/\hbar}$ (reference evolution).
Schr\"odinger states/operators are denoted with subscript $S$,
Heisenberg with $H$, interaction with $I$. Then
\begin{align}
\ket{\psi}_H &\;\equiv\; \op{U}^\dagger(t)\,\ket{\psi(t)}_S
              \;=\; \ket{\psi(0)}_S
\quad(\text{time-independent}), &
\op{O}_H(t) &\;\equiv\; \op{U}^\dagger(t)\,\op{O}_S\,\op{U}(t),
\\[2pt]
\ket{\psi(t)}_I &\;\equiv\; \op{U}_0^\dagger(t)\,\ket{\psi(t)}_S, &
\op{O}_I(t) &\;\equiv\; \op{U}_0^\dagger(t)\,\op{O}_S\,\op{U}_0(t).
\end{align}
The dynamical equations in each picture are
\begin{equation}
\dv{t}\op{O}_H = \tfrac{i}{\hbar}[\op{H},\op{O}_H],
\qquad
i\hbar\,\dv{t}\ket{\psi(t)}_I = \op{V}_I(t)\,\ket{\psi(t)}_I.
\end{equation}
\end{statement}

\begin{remark}[Why the Heisenberg state is constant in time]
The second equality on the Heisenberg line,
$\ket{\psi}_H=\ket{\psi(0)}_S$, is not an additional assumption: it
follows from the Schr\"odinger evolution
$\ket{\psi(t)}_S=\op{U}(t)\ket{\psi(0)}_S$. Substituting into the
definition,
\begin{equation}
\ket{\psi}_H \;=\; \op{U}^\dagger(t)\,\ket{\psi(t)}_S
              \;=\; \op{U}^\dagger(t)\,\op{U}(t)\,\ket{\psi(0)}_S
              \;=\; \ket{\psi(0)}_S.
\end{equation}
The defining unitary $\op{U}^\dagger(t)$ is, by construction, the
inverse of the Schr\"odinger flow, so applying it to
$\ket{\psi(t)}_S$ ``rewinds'' the state back to $t=0$. That is the
sense in which Heisenberg states do not move: all the dynamics has
been absorbed into the operators by the conjugation
$\op{O}_H(t)=\op{U}^\dagger\op{O}_S\op{U}$.
\end{remark}

\begin{proof}[Derivation, interaction picture]
Differentiating $\ket{\psi}_I=\op{U}_0^\dagger\ket{\psi}_S$,
\begin{equation}
i\hbar\,\dv{t}\ket{\psi}_I
= -\op{U}_0^\dagger\op{H}_0\ket{\psi}_S
+ \op{U}_0^\dagger\op{H}\ket{\psi}_S
= \op{U}_0^\dagger\,\op{V}\,\op{U}_0\,\ket{\psi}_I
= \op{V}_I(t)\,\ket{\psi}_I.
\end{equation}
The Heisenberg equation follows analogously by differentiating
$\op{O}_H=\op{U}^\dagger\op{O}_S\op{U}$ and using
$i\hbar\,\dot{\op{U}}=\op{H}\op{U}$.
\end{proof}

\begin{example}[Driven qubit]
For $\op{H}_0=\tfrac{\hbar\omega_q}{2}\hat\sigma_z$ and a drive
$\op{V}(t)=\hbar\Omega\cos(\omega_d t)\,\hat\sigma_x$, moving to the
interaction picture with respect to $\op{H}_0$ rotates
$\hat\sigma_\pm$ at the qubit frequency $\omega_q$, isolating the
near-resonant ($\omega_d\approx\omega_q$) terms. We will use this
exact construction in \cref{sec:rabi,sec:rwa}.
\end{example}

% --------------------------------------------------------------
\section*{Chapter summary and outlook}

We now have everything we need to do physics: states are unit
vectors in a Hilbert space; observables are Hermitian operators;
measurements yield eigenvalues with the probabilities given by the
Born rule; non-commuting observables obey the Robertson uncertainty
relation; and unitary evolution under the Schr\"odinger equation
governs the dynamics. The Heisenberg and interaction pictures give
us complementary, frame-dependent ways of describing the same
physics.

In the next chapter we apply this machinery to the two systems that,
between them, account for almost all of the dynamics encountered in
cQED: the quantum harmonic oscillator (which describes resonators
and, through the Josephson nonlinearity, the leading behaviour of
the transmon) and the two-level system (which describes a qubit and,
in the Bloch-sphere picture, gives the geometric language of
single-qubit gates).
